\section*{4. The calculation of the zeta function of a hypersurface by
Dwork's method}

\subsection*{4.0.}
Let

\begin{equation}\tag{4.0.1}
f(X)=f(X_1,\ldots,X_{n+2})\in\mathbb F_q[X]
\end{equation}

be homogeneous of degree $d$, and let its zero locus define the hypersurface
$V\subset\mathbb P^{n+1}$. Denote by $V^*$ the open subset of $V$ on which
$\prod_{i=1}^{n+2}X_i$ is invertible, and by $V^*_{\mathrm{aff}}$ the affine
cone over $V^*$.

Choose a nontrivial additive character

\begin{equation}\tag{4.0.2}
\chi:\mathbb F_{q^\nu}\longrightarrow\mu_p(C),
\end{equation}

where $C$ is a fixed algebraically closed field of characteristic zero (for
example $C=\mathbb C$). Thus

\begin{equation}\tag{4.0.3}
\text{for }x\in\mathbb F_{q^\nu},\qquad
\sum_{y\in\mathbb F_{q^\nu}}\chi(xy)=
\begin{cases}
q^\nu,&x=0,\\
0,&x\ne0.
\end{cases}
\end{equation}

For $x=(x_1,\ldots,x_{n+2})\in(\mathbb F_{q^\nu}^*)^{n+2}$ we therefore
have

\begin{equation}\tag{4.0.4}
1+\sum_{x_0\in\mathbb F_{q^\nu}^*}\chi(x_0f(x))
=\sum_{x_0\in\mathbb F_{q^\nu}}\chi(x_0f(x))
=
\begin{cases}
q^\nu,&f(x)=0,\\
0,&\text{otherwise}.
\end{cases}
\end{equation}

Summing (4.0.4) over $x\in(\mathbb F_{q^\nu}^*)^{n+2}$ gives

\begin{equation}\tag{4.0.5}
\operatorname{card}V^*_{\mathrm{aff}}(\mathbb F_{q^\nu})
=\frac1{q^\nu}
\sum_{x\in(\mathbb F_{q^\nu}^*)^{n+2}}
\left\{1+\sum_{x_0\in\mathbb F_{q^\nu}^*}\chi(x_0f(x))\right\}.
\end{equation}

Since

\begin{equation}\tag{4.0.6}
\operatorname{card}V^*_{\mathrm{aff}}(\mathbb F_{q^\nu})
=(q^\nu-1)\operatorname{card}V^*(\mathbb F_{q^\nu}),
\end{equation}

formula (4.0.5) becomes

\begin{equation}\tag{4.0.7}
(q^\nu-1)\operatorname{card}V^*(\mathbb F_{q^\nu})
=\frac{(q^\nu-1)^{n+2}}{q^\nu}
+\frac1{q^\nu}
\sum_{(x_0,x)\in(\mathbb F_{q^\nu}^*)^{n+3}}\chi(x_0f(x)).
\end{equation}

Write $f(X)$ explicitly as a sum of monomials,

\begin{equation}\tag{4.0.8}
f(X)=\sum_w a_wX^w,
\end{equation}

where $w=(w_1,\ldots,w_{n+2})$, $\sum_iw_i=d$, and
$X^w\mathrel{\mathop:}=X_1^{w_1}\cdots X_{n+2}^{w_{n+2}}$. Then (4.0.7)
may be rewritten as

\begin{equation}\tag{4.0.9}
(q^\nu-1)\operatorname{card}V^*(\mathbb F_{q^\nu})
=\frac{(q^\nu-1)^{n+2}}{q^\nu}
+\frac1{q^\nu}
\sum_{(x_0,x)\in(\mathbb F_{q^\nu}^*)^{n+3}}
\prod_w\chi(a_wx_0X^w).
\end{equation}

\subsection*{4.1.}
We must now make explicit our ``additive character with values in
characteristic zero.'' One would like simply to take

\begin{equation}\tag{4.1.0}
\chi:x\longmapsto
\exp\!\left(2\pi i\,\operatorname{tr}_{\mathbb F_{q^\nu}/\mathbb F_p}(x)\right)
=\xi_p^{\operatorname{tr}_{\mathbb F_{q^\nu}/\mathbb F_p}(x)}.
\end{equation}

Let

\begin{align}
\tag{4.1.1}
\Omega&=\text{an algebraic closure of }\mathbb Q_p,\\
\tag{4.1.2}
\operatorname{ord}:\Omega^*&\longrightarrow\mathbb Q
  &&\text{be the valuation normalized by }\operatorname{ord}(p)=1,\\
\tag{4.1.3}
\xi_p&\in\Omega
  &&\text{be a primitive $p$-th root of unity.}
\end{align}

One checks (cf. [2]) that there exists

\begin{equation}\tag{4.1.4}
\pi\in\mathbb Q_p(\xi_p),\qquad
\operatorname{ord}(\pi)=\frac1{p-1},
\end{equation}

such that

\begin{equation}\tag{4.1.5}
\sum_{n\geq0}\frac{\pi^{p^n}}{p^n}=0.
\end{equation}

(Notice also that the left-hand side converges in $\mathbb Q_p(\xi_p)$.)

Define the Artin--Hasse exponential

\begin{equation}\tag{4.1.6}
E(X)=\exp\!\left(\sum_{n\geq0}\frac{X^{p^n}}{p^n}\right)
\in\mathbb Q_p[[X]],
\end{equation}

which, miraculously, lies in $\mathbb Z_p[[X]]$. Define

\begin{equation}\tag{4.1.7}
\Theta(X)=E(\pi X).
\end{equation}

Its fundamental properties are summarized in the following lemma.

\subsubsection*{Lemma 4.2.}

\begin{equation}\tag{4.2.1}
\Theta(X)=\sum_n B_n X^n,\qquad B_0=1,\quad B_1=\pi,\quad
\operatorname{ord}(B_n)\geq\frac n{p-1}.
\end{equation}

\begin{equation}\tag{4.2.2}
\Theta(1)=\xi_p,\qquad\text{a primitive $p$-th root of unity.}
\end{equation}

Let $t\in\Omega$, and suppose that, for a given integer $a>0$,

\begin{equation}\tag{4.2.3}
t=t^{p^a};
\end{equation}

that is, $t$ is a Teichmüller representative in
$W(\mathbb F_{p^a})\subset\Omega$. Then
$\sum_{j=0}^{a-1}t^{p^j}\in\mathbb Z_p$, so the formal series

\[
(\Theta(X))^{\sum_{j=0}^{a-1}t^{p^j}}
\]

is defined, and one has the identity of formal power series

\begin{equation}\tag{4.2.3.1}
\prod_{j=0}^{a-1}\Theta(t^{p^j}X)
=(\Theta(X))^{\sum_{j=0}^{a-1}t^{p^j}}.
\end{equation}
